%!TEX root = ../thesis.tex
%*******************************************************************************
%****************************** Fourth Chapter *********************************
%*******************************************************************************

\chapter{Results and Discussion}\label{chap:results-discussion}

% **************************** Define Graphics Path **************************
\ifpdf
    \graphicspath{{Chapter4/Figs/Raster/}{Chapter4/Figs/PDF/}{Chapter4/Figs/}}
\else
    \graphicspath{{Chapter4/Figs/Vector/}{Chapter4/Figs/}}
\fi

\section{Scenario-based Evaluation}
The assistant was evaluated across thirteen multi-turn scenarios designed to mirror the workflows uncovered during requirements elicitation. Each scenario compared the assistant's responses against golden references to assess faithfulness, relevance, and correctness. The evaluation covered medication recall, reference resolution, abbreviation expansion, error handling, disambiguation, free-form reasoning, and patient summarisation tasks in Greek.

\subsection{Scenarios of multi-turn interactions in Greek}
\subsubsection{Happy Paths}
\paragraph{Scenario 1: Memory}
\begin{table}[H]
\centering
\renewcommand{\arraystretch}{1.5} 
\begin{tabular}{|>{\bfseries}l|p{11cm}|}
\hline
Nurse & \textgreek{Η Μαρία Παπαδοπούλου τι αγωγή παίρνει στην παρούσα φάση?} \\
\hline
System - golden & \textgreek{Η Μαρία Παπαδοπούλου παίρνει αλπραζολάμη για τις κρίσεις πανικού με δοσολογία 0,5 mg και πρόγραμμα 2 φορές την ημέρα.} \\
\hline
System - actual & \textgreek{Η Μαρία Παπαδοπούλου λαμβάνει αλπραζολάμη (0.5 m}g) \textgreek{2 φορές την ημέρα για τις κρίσεις πανικού, σύμφωνα με συνταγή του ιατρού Γιάννη Γαστουνιώτη.} \\
\hline
Nurse & \textgreek{Και σε ποια ουσία έχει παρενέργειες?} \\
\hline
System - golden & \textgreek{Στο λίθιο.} \\
\hline
System - actual & \textgreek{Στα έγγραφα της Μαρίας Παπαδοπούλου αναφέρονται παρενέργειες για το λίθιο, το οποίο η ίδια διέκοψε την 1η Οκτωβρίου 2023. Οι παρενέργειες του λιθίου είναι αυξημένη ορμόνη} TSH  \textgreek{(υποθυρεοειδισμός), καταστολή και υπνηλία και τρέμουλο στα χέρια. Δεν αναφέρονται παρενέργειες για την αλπραζολάμη.} \\
\hline
\end{tabular}
\caption{Scenario 1 (audios 1, 2): Memory - MARIA}
\end{table}

\paragraph{Scenario 2: Memory and reference resolution}
\begin{table}[H]
\centering
\renewcommand{\arraystretch}{1.5} 
\begin{tabular}{|>{\bfseries}l|p{11cm}|}
\hline
Nurse & \textgreek{Πότε έγιναν οι προηγούμενες αιματολογικές εξετάσεις του Κώστα Καρά?} \\
\hline
System - golden & \textgreek{Στις 24/11/2024.} \\
\hline
System - actual & \textgreek{Οι προηγούμενες αιματολογικές εξετάσεις του Κώστα Καρά έγιναν στις 10/02/2024 και 24/11/2024.} \\
\hline
Nurse & \textgreek{Και οι επόμενες εξετάσεις του?} \\
\hline
System - golden & \textgreek{Στις 15/5/2025.} \\
\hline
System - actual & \textgreek{Οι επόμενες αιματολογικές εξετάσεις του Κώστα Καρά είναι προγραμματισμένες για τις 15/5/2025.} \\
\hline
Nurse & \textgreek{Τι συστάσεις υπάρχουν?} \\
\hline
System - golden & \textgreek{Προώθηση ψυχοθεραπείας για διαχείριση άγχους και καθοδήγηση για ενίσχυση κοινωνικών δεξιοτήτων.} \\
\hline
System - actual & \textgreek{Οι τρέχουσες ιατρικές συστάσεις για τον Κώστα Καρά περιλαμβάνουν την προώθηση ψυχοθεραπείας για διαχείριση άγχους και καθοδήγηση για ενίσχυση κοινωνικών δεξιοτήτων.} \\
\hline
\end{tabular}
\caption{Scenario 2 (audios 3, 4, 5): Memory and reference resolution - KOSTAS}
\end{table}

\paragraph{Scenario 3: Abbreviations and Acronyms}
\begin{table}[H]
\centering
\renewcommand{\arraystretch}{1.5} 
\begin{tabular}{|>{\bfseries}l|p{11cm}|}
\hline
Nurse & \textgreek{Η Δήμητρα Τσακίρη τι διάγνωση έχει?} \\
\hline
System - golden & \textgreek{ΓΑΔ (διάγνωση το 16/04/2004) και κρίσεις πανικού (διάγνωση το 23/01/2008)} \\
\hline 
System - actual & \textgreek{Η Δήμητρα Τσακίρη έχει διαγνωστεί με ΓΑΔ (Γενικευμένη Αγχώδης Διαταραχή) στις 16/04/2004 και με κρίσεις πανικού στις 23/01/2008.} \\
\hline
Nurse & \textgreek{Πες μου τη δοσολογία της αγωγής για την Ιωάννα Λυμπερίου.} \\
\hline
System - golden & \textgreek{Δραστική ουσία: κουετιαπίνη
Δοσολογία: 300} mg, \textgreek{2 φ/ημ.} \\
\hline
System - actual & \textgreek{Η Ιωάννα Λυμπερίου λαμβάνει κουετιαπίνη σε δοσολογία 300} mg, \textgreek{2 φορές την ημέρα. Η έναρξη της αγωγής ήταν στις 30/08/2018 και συνταγογραφήθηκε από τον καρδιολόγο Ιωάννη Μακρή.} \\
\hline
Nurse & \textgreek{Δώσε μου το τηλέφωνο και το ΑΜΚΑ του Γιώργου Γεωργίου.} \\
\hline
System - golden & \textgreek{Τηλ.: 6974513248 (μητέρα, κα. Αθανασία) και ΑΜΚΑ: 67980831881} \\
\hline
System - actual & \textgreek{Το ΑΜΚΑ του Γιώργου Γεωργίου είναι 67980831881 και το τηλέφωνο είναι 6974513248 (της μητέρας του, κα. Αθανασίας).} \\
\hline
Nurse & \textgreek{Γιώργος Μακρής. Τι διαγνώσεις έχει?} \\
\hline
System - golden & \textgreek{- Κρίσεις πανικού (διάγνωση το 12/11/2002)
- Διπολική διατ. (διάγνωση το 15/01/2021)
- Ψύχωση (διάγνωση το 25/04/2012)} \\
\hline
System - actual & \textgreek{Ο Γιώργος Μακρής έχει διαγνωστεί με τα εξής: Κρίσεις πανικού (12/11/2002), Διπολική διαταραχή (15/01/2021) και Ψύχωση (25/04/2012).} \\
\hline
\end{tabular}
\caption{Scenario 3 (audios 6, 7, 8, 9): Abbreviations and Acronyms - PATIENT 01 and PATIENT 03 and PATIENT 04 and PATIENT 05}
\end{table}

 
\subsubsection{Error Scenarios}
\paragraph{Scenario 4: Non-existent patient}
\begin{table}[H]
\centering
\renewcommand{\arraystretch}{1.5} 
\begin{tabular}{|>{\bfseries}l|p{11cm}|}
\hline
Nurse & \textgreek{Σε ποια ουσία έχει παρενέργειες η Δήμητρα Καρακίτσου?} \\
\hline
System - golden & \textgreek{Δεν ξέρω.} \\
\hline
System - actual & \textgreek{Δεν βρέθηκαν πληροφορίες σχετικά με το σε ποια ουσία έχει παρενέργειες η Δήμητρα Καρακίτσου.} \\
\hline
\end{tabular}
\caption{Scenario 4 (audio 10): Non-existent patient}
\end{table}

\paragraph{Scenario 5: Non-existent patient who shares the same last name with a known patient}
\begin{table}[H]
\centering
\renewcommand{\arraystretch}{1.5} 
\begin{tabular}{|>{\bfseries}l|p{11cm}|}
\hline
Nurse & \textgreek{Πες μου τις συστάσεις για την ασθενή Θάλεια Παπαδοπούλου.} \\
\hline
System - golden & \textgreek{Δεν ξέρω.} \\
\hline
System - actual & \textgreek{Δεν βρέθηκαν πληροφορίες για την ασθενή Θάλεια Παπαδοπούλου.} \\
\hline
\end{tabular}
\caption{Scenario 5 (audio 11): Non-existent patient who shares the same last name with a known patient}
\end{table}

\paragraph{Scenario 6: Known patient with non-existent info}
\begin{table}[H]
\centering
\renewcommand{\arraystretch}{1.5} 
\begin{tabular}{|>{\bfseries}l|p{11cm}|}
\hline
Nurse & \textgreek{Πότε είναι η επόμενη ενέσιμη αγωγή του Κώστα Καρά?} \\
\hline
System - golden & \textgreek{Δεν έχει ενέσιμη αγωγή.} \\
\hline
System - actual & \textgreek{Ο Κώστας Καράς δεν έχει ενέσιμη αγωγή. Έχει όμως προγραμματισμένες αιματολογικές εξετάσεις στις 15/5/2025.} \\
\hline
\end{tabular}
\caption{Scenario 6 (audio 12): Known patient with non-existent info - KOSTAS}
\end{table}

\paragraph{Scenario 7: Doctor's name}
\begin{table}[H]
\centering
\renewcommand{\arraystretch}{1.5} 
\begin{tabular}{|>{\bfseries}l|p{11cm}|}
\hline
Nurse & \textgreek{Πες μου το ΑΜΚΑ της Κωνσταντίνας Λυμπερίου.} \\
\hline
System - golden & \textgreek{Δεν ξέρω.} \\
\hline
System - actual & \textgreek{Δεν βρέθηκαν δεδομένα για την Κωνσταντίνα Λυμπερίου. Δεν μπορώ να σου δώσω το ΑΜΚΑ της.} \\
\hline
\end{tabular}
\caption{Scenario 7 (audio 13): Doctor's name - GEORGIA M}
\end{table}

\subsubsection{Disambiguation Scenarios}
\paragraph{Scenario 8: Disambiguation of identical patient name}
\begin{table}[H]
\centering
\renewcommand{\arraystretch}{1.5} 
\begin{tabular}{|>{\bfseries}l|p{11cm}|}
\hline
Nurse & \textgreek{Θύμισέ μου το κοινωνικό ιστορικό της Ελένης Κανελλοπούλου.} \\
\hline
System & \textgreek{Έχω βρει πληροφορίες για δύο ασθενείς με το όνομα Ελένη Κανελλοπούλου. Η μία είναι 54 ετών (ΑΜΚΑ: 63681990357) και η άλλη είναι 41 ετών (ΑΜΚΑ: 65357602736). Για ποια από τις δύο Ελένες Κανελλοπούλου θέλετε το κοινωνικό ιστορικό?} \\
\hline
Nurse & \textgreek{Πες μου για τη μεγαλύτερη.} \\
\hline
System - golden & PATIENT 15.txt \\
\hline
System - actual & \textgreek{OK. Το κοινωνικό ιστορικό της Ελένης Κανελλοπούλου (ΑΜΚΑ: 63681990357) είναι:
*   Μένει μόνη σε ενοικιαζόμενο διαμέρισμα
*   Έχει ελάχιστη οικογενειακή υποστήριξη
*   Δεν εργάζεται επί του παρόντος
*   Πρώην υπάλληλος γραφείου, αποχώρησε λόγω επιδείνωσης της ψυχικής υγείας
*   Διατηρεί περιορισμένη επαφή με μία εξαδέλφη. Όσον αφορά τις συνήθειες και την καθημερινότητά της:
*   Καπνίζει περίπου 10 τσιγάρα την ημέρα
*   Δεν καταναλώνει αλκοόλ
*   Παρουσιάζει διαταραχές ύπνου (δυσκολία έναρξης και διατήρησης ύπνου, συχνές διακοπές για ούρηση)
*   Βλέπει τηλεόραση και ράβει για να ηρεμεί.
Οι συστάσεις για την Ελένη Κανελλοπούλου είναι:
*   Συστηματική ψυχοθεραπεία και παρακολούθηση
*   Ενίσχυση κοινωνικών δεξιοτήτων μέσω ομάδων υποστήριξης
*   Εκπαίδευση στη διαχείριση στρες και τεχνικές αναπνοής
*   Συμμετοχή σε πρόγραμμα επανένταξης στην εργασία (ΔΥΠΑ)
*   Προγραμματισμός αξιολόγησης από Κοινωνική Υπηρεσία για πιθανή παροχή βοήθειας κατ’ οίκον.} \\
\hline
\end{tabular}
\caption{Scenario 8 (audio 14, 15): Disambiguation of identical patient name - PATIENT 13 and PATIENT 15}
\end{table}

\paragraph{Scenario 9: Disambiguation of identical patient name and double intent}
\begin{table}[H]
\centering
\renewcommand{\arraystretch}{1.5} 
\begin{tabular}{|>{\bfseries}l|p{11cm}|}
\hline
Nurse & \textgreek{Πότε νοσηλεύτηκε τελευταία φορά η Δήμητρα Μακρή και τι αγωγή παίρνει?} \\
\hline
System & \textgreek{Υπάρχουν δύο ασθενείς με το όνομα Δήμητρα Μακρή. Για να σας δώσω τις σωστές πληροφορίες, θα μπορούσατε να επιβεβαιώσετε τον ΑΜΚΑ ή την ηλικία της ασθενούς για την οποία ενδιαφέρεστε?} \\
\hline
Nurse & \textgreek{Δε θυμάμαι το ΑΜΚΑ της. Πες μου για την πιο μικρή.} \\
\hline
System - golden &  \textgreek{Στις 20/03/2013: Επιθετικό-μανιακό πεισόδιο με παραμονή 2 ημέρες, και στις 27/09/2017: Επεισόδιο σχιζοφρένειας με παραμονή 1 ημέρες. Παίρνει ρισπεριδόνη, 3 mg, 2 φορές την ημέρα, και ζιπραζιδόνη, 40 mg, 2 φορές την ημέρα μαζί με φαγητό.} \\
\hline
System - actual & \textgreek{Η Δήμητρα Μακρή, 32 ετών (ΑΜΚΑ 41661927592), νοσηλεύτηκε τελευταία φορά στις 27/09/2017 και λαμβάνει ρισπεριδόνη (3 mg, 2 φ. την ημέρα) και ζιπραζιδόνη (40 mg, 2 φ. την ημέρα μαζί με φαγητό).} \\
\hline
\end{tabular}
\caption{Scenario 9 (audio 16, 17): Disambiguation of identical patient name and double intent - PATIENT 10 and PATIENT 16}
\end{table}

\paragraph{Scenario 10: Free-form text}
\begin{table}[H]
\centering
\renewcommand{\arraystretch}{1.5} 
\begin{tabular}{|>{\bfseries}l|p{11cm}|}
\hline
Nurse & \textgreek{Θύμισέ μου αν ακολουθεί τη θεραπεία της η Βασιλική Αντωνίου.} \\
\hline
System - golden & \textgreek{Λαμβάνει τα φάρμακά της με σχετική συνέπεια, αν και παρατηρούνται περιστασιακές παραλείψεις, κυρίως κατά τις περιόδους ύφεσης. 
Παρουσιάζει μειωμένη αυτογνωσία σχετικά με την ανάγκη για συνεχή θεραπεία, γεγονός που αποδίδεται εν μέρει σε μερική ανοσογνωσία.} \\
\hline
System - actual & \textgreek{Η Βασιλική Αντωνίου λαμβάνει τα φάρμακά της με σχετική συνέπεια, αν και παρατηρούνται περιστασιακές παραλείψεις, κυρίως κατά τις περιόδους ύφεσης. Παρουσιάζει μειωμένη αυτογνωσία σχετικά με την ανάγκη για συνεχή θεραπεία, γεγονός που αποδίδεται εν μέρει σε μερική ανοσογνωσία.} \\
\hline
Nurse & \textgreek{Έχει κάποιες συστάσεις στο φάκελό της?} \\
\hline
System - golden & \textgreek{Οι θεραπευτικές κατευθύνσεις περιλαμβάνουν τη συνέχιση της φαρμακευτικής αγωγής υπό τακτική παρακολούθηση, την επανένταξή της σε ομάδα υποστηρικτικής ψυχοθεραπείας σε εβδομαδιαία βάση, την καθιέρωση ενός σταθερού προγράμματος δραστηριοτήτων που να περιλαμβάνει βόλτες και ελαφριές οικιακές εργασίες, καθώς και τη συμμετοχή σε πρόγραμμα κοινωνικοποίησης στο κέντρο ημέρας. 
Επιπλέον, προτείνεται παραπομπή σε ψυχολόγο για γνωσιακή-συμπεριφορική θεραπεία, παροχή συμβουλευτικής υποστήριξης προς την οικογένειά της, ενίσχυση της αίσθησης αυτοαποτελεσματικότητας και ενδυνάμωσης, και μηνιαία παρακολούθηση από κοινωνικό λειτουργό.} \\
\hline
System - actual & \textgreek{Ναι, υπάρχουν κάποιες συστάσεις στο φάκελο της Βασιλικής Αντωνίου. Προτείνεται παραπομπή σε ψυχολόγο για γνωσιακή-συμπεριφορική θεραπεία, παροχή συμβουλευτικής υποστήριξης προς την οικογένειά της, ενίσχυση της αίσθησης αυτοαποτελεσματικότητας και ενδυνάμωσης, και μηνιαία παρακολούθηση από κοινωνικό λειτουργό. Επιπλέον, οι θεραπευτικές κατευθύνσεις περιλαμβάνουν τη συνέχιση της φαρμακευτικής αγωγής υπό τακτική παρακολούθηση, την επανένταξή της σε ομάδα υποστηρικτικής ψυχοθεραπείας σε εβδομαδιαία βάση, την καθιέρωση ενός σταθερού προγράμματος δραστηριοτήτων που να περιλαμβάνει βόλτες και ελαφριές οικιακές εργασίες, καθώς και τη συμμετοχή σε πρόγραμμα κοινωνικοποίησης στο κέντρο ημέρας.} \\
\hline
\end{tabular}
\caption{Scenario 10 (audio 18, 19): Free-form text - PATIENT 40}
\end{table}

\paragraph{Scenario 11: Free-form text}
\begin{table}[H]
\centering
\renewcommand{\arraystretch}{1.5} 
\begin{tabular}{|>{\bfseries}l|p{11cm}|}
\hline
Nurse & \textgreek{Ποια είναι η τρέχουσα ψυχολογική κατάσταση του Νίκου Τρανού?} \\
\hline
System - golden & \textgreek{Κατά την τελευταία ψυχολογική εκτίμηση, αναφέρθηκαν εναλλαγές διάθεσης με περιόδους απογοήτευσης και μειωμένο ενδιαφέρον για δραστηριότητες, χωρίς ωστόσο ενδείξεις ενεργού αυτοκτονικού ιδεασμού.} \\
\hline
System - actual & \textgreek{Η τρέχουσα ψυχολογική κατάσταση του Νίκου Τρανού (ΑΜΚΑ: 71462734135), σύμφωνα με την ψυχολογική αξιολόγηση, παρουσιάζει εναλλαγές στη διάθεση και περιστασιακή απογοήτευση. Η ενέργειά του είναι μέτρια, με μειωμένη όρεξη για δραστηριότητες, ενώ η κοινωνική του συμπεριφορά είναι περιορισμένη. Δεν έχουν αναφερθεί ενδείξεις αυτοκτονικού ιδεασμού.} \\
\hline
Nurse & \textgreek{Ακολουθεί την αγωγή του?} \\
\hline
System - golden & \textgreek{Η συμμόρφωση του ασθενούς στη θεραπεία είναι μέτρια. Ο ίδιος αναφέρει ότι λαμβάνει τα φάρμακά του με συνέπεια, ωστόσο παραδέχεται πως τα παραλείπει περιστασιακά, ιδίως όταν αισθάνεται πως «είναι καλά» — φαινόμενο συχνό σε χρόνιους ψυχιατρικούς ασθενείς.} \\
\hline
System - actual & \textgreek{Ο Νίκος Τρανός (ΑΜΚΑ: 71462734135) παίρνει τα φάρμακά του με σχετική συνέπεια, αλλά αναφέρει περιστασιακές παραλείψεις όταν νιώθει καλά.} \\
\hline
\end{tabular}
\caption{Scenario 11 (audio 20, 21): Free-form text - PATIENT 38}
\end{table}

\begin{table}[H]
\centering
\renewcommand{\arraystretch}{1.5} 
\begin{tabular}{|>{\bfseries}l|p{11cm}|}
\hline
Nurse & \textgreek{Δείχνει επιθυμία για απομόνωση? Πώς είναι κοινωνικά?} \\
\hline
System - golden & \textgreek{Η κοινωνική του αλληλεπίδραση παραμένει περιορισμένη, ενώ φαίνεται να βιώνει συναισθήματα απομόνωσης, κάτι που επιτείνεται από την ελλιπή οικογενειακή στήριξη. Διαμένει μόνος του σε διαμέρισμα που του παραχωρήθηκε από την κοινότητα, ενώ η επικοινωνία με την οικογένειά του είναι περιορισμένη — διατηρεί επαφή μόνο με έναν αδερφό του, τον οποίο βλέπει περιστασιακά. Δεν έχει δημιουργήσει σταθερό οικογενειακό κύκλο ή στενές φιλίες, αν και τον τελευταίο χρόνο συμμετέχει σποραδικά σε δραστηριότητες Κέντρου Ημέρας για άτομα με ψυχικές διαταραχές, όπου φαίνεται να δείχνει διστακτικό αλλά θετικό ενδιαφέρον για κοινωνική αλληλεπίδραση. Ανέφερε επίσης ότι κάποιες Κυριακές πηγαίνει σε καφενείο της γειτονιάς, κυρίως για να «παρακολουθεί τους άλλους» και όχι για να συμμετάσχει ενεργά στη συζήτηση.} \\
\hline
System - actual & \textgreek{Ο Νίκος Τρανός έχει περιορισμένη κοινωνική αλληλεπίδραση, γεγονός που υποδηλώνει επιθυμία για απομόνωση.} \\
\hline
Nurse & \textgreek{Απαιτείται ενημέρωση ιατρού ή παρέμβαση ψυχολόγου ή ψυχιάτρου?} \\
\hline
System - golden & \textgreek{Ο συνολικός ψυχιατρικός του φάκελος υποδεικνύει την ανάγκη για πιο στενή ψυχιατρική παρακολούθηση, ενίσχυση της ψυχοκοινωνικής υποστήριξης, και πιθανή ένταξή του σε προγράμματα ψυχοεκπαίδευσης και ενδυνάμωσης, ώστε να βελτιωθεί η ποιότητα ζωής και να περιοριστούν οι υποτροπές.} \\
\hline
System - actual & \textgreek{Με βάση τα έγγραφα του ασθενούς Νίκου Τρανού, λαμβάνοντας υπόψη τις διαγνώσεις του (ψύχωση, μανιοκατάθλιψη), το ιστορικό επεισοδίων, την τρέχουσα φαρμακευτική αγωγή (ολανζαπίνη), τις παρενέργειες που είχε από προηγούμενη φαρμακευτική αγωγή (φλουοξετίνη), την περιορισμένη κοινωνική συμπεριφορά και τις εναλλαγές στη διάθεση, είναι απαραίτητη η ενημέρωση του θεράποντος ιατρού του, Ιωάννη Μακρή, καθώς και η ψυχολογική υποστήριξη από ψυχολόγο ή ψυχίατρο.} \\
\hline
\end{tabular}
\caption{Scenario 11 (audio 22, 23): Free-form text - PATIENT 38}
\end{table}

\subsubsection{Summarization Scenarios}
\paragraph{Scenario 12: Summary of a patient's record from structured text}
\begin{table}[H]
\centering
\renewcommand{\arraystretch}{1.5} 
\begin{tabular}{|>{\bfseries}l|p{11cm}|}
\hline
Nurse & \textgreek{Δώσε μου τη γενική εικόνα για τον Θανάση Πράσινο.} \\
\hline
System - golden & PATIENT 29.txt \\
\hline
System - actual & \textgreek{Η γενική εικόνα για τον Θανάση Πράσινο είναι η εξής: Ο Θανάσης Πράσινος, 24 ετών, με ΑΜΚΑ 54565609013, έχει διαγνωστεί με μανιοκατάθλιψη (από 30/04/2014) και σχιζοφρένεια (από 11/01/2005). Είχε ένα επεισόδιο στις 25/03/2024 που χρειάστηκε νοσηλεία για 1 ημέρα.
Αυτή τη στιγμή λαμβάνει κουετιαπίνη 300mg, δύο φορές την ημέρα, για την αντιμετώπιση των συμπτωμάτων του, με συνταγογράφο τον ψυχολόγο Δημήτριο Κατσαρό (έναρξη 25/03/2024). Στο παρελθόν είχε παρενέργειες (ναυτία και μείωση λίμπιντο) από την ολανζαπίνη, η οποία διακόπηκε στις 23/05/2024.
Η διάθεσή του παρουσιάζει απότομες εναλλαγές και έχει μέτρια ενέργεια, με μειωμένη όρεξη για δραστηριότητες. Έχει παρατηρηθεί περίεργη/βίαιη συμπεριφορά στα} social media (TikTok, Facebook) \textgreek{σε περιόδους έξαρσης της ασθένειας. Δεν έχουν αναφερθεί ενδείξεις αυτοκτονικού ιδεασμού.
Η οικογένειά του αποφεύγει να συζητάει για το πρόβλημά του. Εργάζεται ως γυμναστής στα} Alter Life.
\textgreek{Συνιστάται η συνέχιση της συστηματικής ψυχοθεραπείας και παρακολούθησης. Παίρνει τα φάρμακά του με σχετική συνέπεια, αλλά αναφέρει περιστασιακές παραλείψεις όταν νιώθει καλά. Έχει κάνει αιματολογικές εξετάσεις στις 31/05/2024 και 23/11/2024, και οι επόμενες είναι προγραμματισμένες για τις 06/10/2025.} \\
\hline
\end{tabular}
\caption{Scenario 12 (audio 24): Summary of a patient's record from structured text - PATIENT 29}
\end{table}

\paragraph{Scenario 13: Summary of a patient's record from free-form text}
\begin{table}[H]
\centering
\renewcommand{\arraystretch}{1.5} 
\begin{tabular}{|>{\bfseries}l|p{11cm}|}
\hline
Nurse & \textgreek{Πες μου κάποια σημαντικά πράγματα που πρέπει να ξέρω για τη Χριστίνα Παγούτση.} \\
\hline
System - golden & PATIENT 41.txt \\
\hline
System - actual & \textgreek{Εδώ είναι μερικά σημαντικά πράγματα που πρέπει να γνωρίζετε για την Χριστίνα Παγούτση:
**Ιστορικό:**
*   Διαγνώστηκε με ιδεοψυχαναγκαστική διαταραχή (από 11 Ιουνίου 2010), μείζονα καταθλιπτική διαταραχή (από 25 Ιανουαρίου 2006), παραληρητική ψύχωση (από 1 Δεκεμβρίου 2005) και αϋπνία σχετιζόμενη με αγχώδη διαταραχή (από το 2021).
*   Νοσηλεύτηκε στις 4 Οκτωβρίου 2008 για τέσσερις ημέρες λόγω έντονων εμμονικών σκέψεων με αυτοτραυματικές τάσεις.
*   Στις 12 Σεπτεμβρίου 2020 υπήρξε παρέμβαση κινητής μονάδας, εξαιτίας επιδείνωσης των συμπτωμάτων.
*   Είχε ένα ψυχωσικό επεισόδιο στις 21 Απριλίου 2013, που οδήγησε σε νοσηλεία δύο ημερών.
*   **Φαρμακευτική αγωγή:**
*   Λαμβάνει σερτραλίνη 100 mg ημερησίως, ρισπεριδόνη 2 mg ημερησίως, τραζοδόνη 50 mg πριν τον ύπνο και λοραζεπάμη 1 mg κατά περίπτωση.
*   Στο παρελθόν λάμβανε αλπραζολάμη, η οποία διακόπηκε στις 6 Μαρτίου 2023.
*   Έχει μέτρια συμμόρφωση στη φαρμακευτική αγωγή και ιστορικό διακοπής της αγωγής χωρίς συνεννόηση με το θεραπευτικό προσωπικό.
*   **Εργαστηριακές εξετάσεις:**
*   Πρόσφατες αιματολογικές εξετάσεις (17 Αυγούστου 2024) έδειξαν μέτρια αυξημένα τριγλυκερίδια.
*   Στις 8 Ιουλίου 2023 είχε εμφανίσει ελαφρά αναιμία και χαμηλή φερριτίνη. Η επόμενη εξέταση έχει προγραμματιστεί για τις 29 Δεκεμβρίου 2025.
*   **Κοινωνικό προφίλ:**
*   Διαμένει μόνη της σε ενοικιαζόμενο διαμέρισμα με περιορισμένους οικονομικούς πόρους.
*   **Δημογραφικά στοιχεία:**
*   Είναι 59 ετών. ΑΜΚΑ 93599718720.} \\
\hline
\end{tabular}
\caption{Scenario 13 (audio 25): Summary of a patient's record from free-form text - PATIENT 41}
\end{table}




\section{Interpretation in Relation to the Research Questions}
The results demonstrate the feasibility of deploying a retrieval-augmented conversational assistant for psychiatric nurses (RQ1). Across the happy-path scenarios, the assistant reproduced medication plans, laboratory schedules, and clinician recommendations with high fidelity, even when disambiguation or reference resolution was required. Performance on the error scenarios indicates that the system can decline to answer when information is missing, reducing the risk of unsupported claims. The detailed free-form and summarisation tasks show that the assistant can synthesise longer passages while remaining grounded in the retrieved notes, supporting positive answers to RQ2. RQ3 remains partially answered: while users reported enthusiasm during interviews, full validation of experience and satisfaction requires real-world deployments with practicing nurses.

\section{Comparison with Literature}
The observed results align with the literature on RAG-enabled healthcare assistants that emphasise improved information extraction and transparency \cite{clinicalentityIE2025}. The assistant's ability to ground responses in curated notes addresses concerns about hallucinations raised in Section~\ref{hallucsLLM}, while the accurate handling of disambiguation mirrors best practices found in prior chronic disease and mental health agents \cite{max2021, schizophrenia2010}. Unlike general-purpose assistants that struggle with medical terminology \cite{healthprompts2020, mouth2019}, the domain-specific retrieval and prompt engineering strategies yielded consistent terminology expansion and context tracking.

\section{Implications and Contributions}
Practically, the evaluation suggests that nurses could rely on the assistant to retrieve patient-specific information rapidly, potentially reducing time spent consulting handwritten records. Theoretical contributions include an applied demonstration of agentic RAG workflows for Greek-language clinical content, expanding on existing literature that has largely focused on English-speaking contexts. The results also highlight the value of combining structured and free-form notes within a single retrieval corpus, illustrating a pathway for future nurse-facing tools.

\section{Limitations}
While the assistant performed reliably on synthetic data, the absence of real patient records and clinical deployment means that generalisability is not yet proven. Certain free-form answers remained more verbose than the golden references, indicating the need for tighter prompt constraints to ensure succinctness. Additionally, automated metrics such as RAGAS were only moderately strong, underscoring the importance of human evaluation alongside quantitative scores. These limitations echo the responsible AI considerations discussed in Chapter~\ref{chap:methodology} and will guide future iterations of the system.

